% Options for packages loaded elsewhere
% Options for packages loaded elsewhere
\PassOptionsToPackage{unicode}{hyperref}
\PassOptionsToPackage{hyphens}{url}
\PassOptionsToPackage{dvipsnames,svgnames,x11names}{xcolor}
%
\documentclass[
  letterpaper,
  DIV=11,
  numbers=noendperiod]{scrartcl}
\usepackage{xcolor}
\usepackage{amsmath,amssymb}
\setcounter{secnumdepth}{-\maxdimen} % remove section numbering
\usepackage{iftex}
\ifPDFTeX
  \usepackage[T1]{fontenc}
  \usepackage[utf8]{inputenc}
  \usepackage{textcomp} % provide euro and other symbols
\else % if luatex or xetex
  \usepackage{unicode-math} % this also loads fontspec
  \defaultfontfeatures{Scale=MatchLowercase}
  \defaultfontfeatures[\rmfamily]{Ligatures=TeX,Scale=1}
\fi
\usepackage{lmodern}
\ifPDFTeX\else
  % xetex/luatex font selection
\fi
% Use upquote if available, for straight quotes in verbatim environments
\IfFileExists{upquote.sty}{\usepackage{upquote}}{}
\IfFileExists{microtype.sty}{% use microtype if available
  \usepackage[]{microtype}
  \UseMicrotypeSet[protrusion]{basicmath} % disable protrusion for tt fonts
}{}
\makeatletter
\@ifundefined{KOMAClassName}{% if non-KOMA class
  \IfFileExists{parskip.sty}{%
    \usepackage{parskip}
  }{% else
    \setlength{\parindent}{0pt}
    \setlength{\parskip}{6pt plus 2pt minus 1pt}}
}{% if KOMA class
  \KOMAoptions{parskip=half}}
\makeatother
% Make \paragraph and \subparagraph free-standing
\makeatletter
\ifx\paragraph\undefined\else
  \let\oldparagraph\paragraph
  \renewcommand{\paragraph}{
    \@ifstar
      \xxxParagraphStar
      \xxxParagraphNoStar
  }
  \newcommand{\xxxParagraphStar}[1]{\oldparagraph*{#1}\mbox{}}
  \newcommand{\xxxParagraphNoStar}[1]{\oldparagraph{#1}\mbox{}}
\fi
\ifx\subparagraph\undefined\else
  \let\oldsubparagraph\subparagraph
  \renewcommand{\subparagraph}{
    \@ifstar
      \xxxSubParagraphStar
      \xxxSubParagraphNoStar
  }
  \newcommand{\xxxSubParagraphStar}[1]{\oldsubparagraph*{#1}\mbox{}}
  \newcommand{\xxxSubParagraphNoStar}[1]{\oldsubparagraph{#1}\mbox{}}
\fi
\makeatother

\usepackage{color}
\usepackage{fancyvrb}
\newcommand{\VerbBar}{|}
\newcommand{\VERB}{\Verb[commandchars=\\\{\}]}
\DefineVerbatimEnvironment{Highlighting}{Verbatim}{commandchars=\\\{\}}
% Add ',fontsize=\small' for more characters per line
\usepackage{framed}
\definecolor{shadecolor}{RGB}{241,243,245}
\newenvironment{Shaded}{\begin{snugshade}}{\end{snugshade}}
\newcommand{\AlertTok}[1]{\textcolor[rgb]{0.68,0.00,0.00}{#1}}
\newcommand{\AnnotationTok}[1]{\textcolor[rgb]{0.37,0.37,0.37}{#1}}
\newcommand{\AttributeTok}[1]{\textcolor[rgb]{0.40,0.45,0.13}{#1}}
\newcommand{\BaseNTok}[1]{\textcolor[rgb]{0.68,0.00,0.00}{#1}}
\newcommand{\BuiltInTok}[1]{\textcolor[rgb]{0.00,0.23,0.31}{#1}}
\newcommand{\CharTok}[1]{\textcolor[rgb]{0.13,0.47,0.30}{#1}}
\newcommand{\CommentTok}[1]{\textcolor[rgb]{0.37,0.37,0.37}{#1}}
\newcommand{\CommentVarTok}[1]{\textcolor[rgb]{0.37,0.37,0.37}{\textit{#1}}}
\newcommand{\ConstantTok}[1]{\textcolor[rgb]{0.56,0.35,0.01}{#1}}
\newcommand{\ControlFlowTok}[1]{\textcolor[rgb]{0.00,0.23,0.31}{\textbf{#1}}}
\newcommand{\DataTypeTok}[1]{\textcolor[rgb]{0.68,0.00,0.00}{#1}}
\newcommand{\DecValTok}[1]{\textcolor[rgb]{0.68,0.00,0.00}{#1}}
\newcommand{\DocumentationTok}[1]{\textcolor[rgb]{0.37,0.37,0.37}{\textit{#1}}}
\newcommand{\ErrorTok}[1]{\textcolor[rgb]{0.68,0.00,0.00}{#1}}
\newcommand{\ExtensionTok}[1]{\textcolor[rgb]{0.00,0.23,0.31}{#1}}
\newcommand{\FloatTok}[1]{\textcolor[rgb]{0.68,0.00,0.00}{#1}}
\newcommand{\FunctionTok}[1]{\textcolor[rgb]{0.28,0.35,0.67}{#1}}
\newcommand{\ImportTok}[1]{\textcolor[rgb]{0.00,0.46,0.62}{#1}}
\newcommand{\InformationTok}[1]{\textcolor[rgb]{0.37,0.37,0.37}{#1}}
\newcommand{\KeywordTok}[1]{\textcolor[rgb]{0.00,0.23,0.31}{\textbf{#1}}}
\newcommand{\NormalTok}[1]{\textcolor[rgb]{0.00,0.23,0.31}{#1}}
\newcommand{\OperatorTok}[1]{\textcolor[rgb]{0.37,0.37,0.37}{#1}}
\newcommand{\OtherTok}[1]{\textcolor[rgb]{0.00,0.23,0.31}{#1}}
\newcommand{\PreprocessorTok}[1]{\textcolor[rgb]{0.68,0.00,0.00}{#1}}
\newcommand{\RegionMarkerTok}[1]{\textcolor[rgb]{0.00,0.23,0.31}{#1}}
\newcommand{\SpecialCharTok}[1]{\textcolor[rgb]{0.37,0.37,0.37}{#1}}
\newcommand{\SpecialStringTok}[1]{\textcolor[rgb]{0.13,0.47,0.30}{#1}}
\newcommand{\StringTok}[1]{\textcolor[rgb]{0.13,0.47,0.30}{#1}}
\newcommand{\VariableTok}[1]{\textcolor[rgb]{0.07,0.07,0.07}{#1}}
\newcommand{\VerbatimStringTok}[1]{\textcolor[rgb]{0.13,0.47,0.30}{#1}}
\newcommand{\WarningTok}[1]{\textcolor[rgb]{0.37,0.37,0.37}{\textit{#1}}}

\usepackage{longtable,booktabs,array}
\newcounter{none} % for unnumbered tables
\usepackage{calc} % for calculating minipage widths
% Correct order of tables after \paragraph or \subparagraph
\usepackage{etoolbox}
\makeatletter
\patchcmd\longtable{\par}{\if@noskipsec\mbox{}\fi\par}{}{}
\makeatother
% Allow footnotes in longtable head/foot
\IfFileExists{footnotehyper.sty}{\usepackage{footnotehyper}}{\usepackage{footnote}}
\makesavenoteenv{longtable}
\usepackage{graphicx}
\makeatletter
\newsavebox\pandoc@box
\newcommand*\pandocbounded[1]{% scales image to fit in text height/width
  \sbox\pandoc@box{#1}%
  \Gscale@div\@tempa{\textheight}{\dimexpr\ht\pandoc@box+\dp\pandoc@box\relax}%
  \Gscale@div\@tempb{\linewidth}{\wd\pandoc@box}%
  \ifdim\@tempb\p@<\@tempa\p@\let\@tempa\@tempb\fi% select the smaller of both
  \ifdim\@tempa\p@<\p@\scalebox{\@tempa}{\usebox\pandoc@box}%
  \else\usebox{\pandoc@box}%
  \fi%
}
% Set default figure placement to htbp
\def\fps@figure{htbp}
\makeatother





\setlength{\emergencystretch}{3em} % prevent overfull lines

\providecommand{\tightlist}{%
  \setlength{\itemsep}{0pt}\setlength{\parskip}{0pt}}



 


\KOMAoption{captions}{tableheading}
\makeatletter
\@ifpackageloaded{caption}{}{\usepackage{caption}}
\AtBeginDocument{%
\ifdefined\contentsname
  \renewcommand*\contentsname{Table of contents}
\else
  \newcommand\contentsname{Table of contents}
\fi
\ifdefined\listfigurename
  \renewcommand*\listfigurename{List of Figures}
\else
  \newcommand\listfigurename{List of Figures}
\fi
\ifdefined\listtablename
  \renewcommand*\listtablename{List of Tables}
\else
  \newcommand\listtablename{List of Tables}
\fi
\ifdefined\figurename
  \renewcommand*\figurename{Figure}
\else
  \newcommand\figurename{Figure}
\fi
\ifdefined\tablename
  \renewcommand*\tablename{Table}
\else
  \newcommand\tablename{Table}
\fi
}
\@ifpackageloaded{float}{}{\usepackage{float}}
\floatstyle{ruled}
\@ifundefined{c@chapter}{\newfloat{codelisting}{h}{lop}}{\newfloat{codelisting}{h}{lop}[chapter]}
\floatname{codelisting}{Listing}
\newcommand*\listoflistings{\listof{codelisting}{List of Listings}}
\makeatother
\makeatletter
\makeatother
\makeatletter
\@ifpackageloaded{caption}{}{\usepackage{caption}}
\@ifpackageloaded{subcaption}{}{\usepackage{subcaption}}
\makeatother
\usepackage{bookmark}
\IfFileExists{xurl.sty}{\usepackage{xurl}}{} % add URL line breaks if available
\urlstyle{same}
\hypersetup{
  pdftitle={Practical 6: Prediction using GWAS summary statistics},
  pdfauthor={Summer Institute of Statical Genetics (Module 15)},
  colorlinks=true,
  linkcolor={blue},
  filecolor={Maroon},
  citecolor={Blue},
  urlcolor={Blue},
  pdfcreator={LaTeX via pandoc}}


\title{Practical 6: Prediction using GWAS summary statistics}
\author{Summer Institute of Statical Genetics (Module 15)}
\date{2023-07-25}
\begin{document}
\maketitle


{\textbf{This practical runs on the SISG2023 server.}}

\subsection{Preparation}\label{preparation}

This practical proposes a detailed tutorial for running three popular
polygenic score (PGS) methods using GWAS summary statistics as
introduced in \textbf{Lecture 6}.

\textbf{Method 1: C+PT}

The first method is classically denoted as ``Clumping + P-value
Thresholding (C+PT)''. This method is also abbreviated as P+T or C+T in
certain publications. In brief, the principle of this method is compare
various sets of uncorrelated SNPs (e.g., maximum squared pairwise
correlation between allele counts at SNPs in the selected set is
\(r^2 \leq 0.1\)) that are associated with the trait / disease as
certain p-value threshold (e.g., p\textless0.01). We then select the set
of SNPs that yields the largest prediction accuracy with the trait /
disease of interest in a validation set. This method is broadly used
because of it simplicity but may not often yield the largest accuracy.
In this practical, we will use \texttt{PLINK} and \texttt{R} to
determine these optimal sets of SNPs. PGS will be calculated using
marginal/GWAS SNP effects as weights.

\textbf{Method 2: SBLUP}

The second method, the Summary-data based BLUP, is an approximation of
the BLUP prediction method introduced in \textbf{Lecture 5}. In brief,
SBLUP calculates PGS weights \(\mathbf{b}\) using the following equation

\[ \mathbf{b} = \left( \mathbf{R} + \frac{M_{SNP}(1-h_{SNP}^2)}{N_{GWAS}h_{SNP}^2}\mathbf{I}_M \right)^{-1} \boldsymbol{\beta} \]
where \(\boldsymbol{\beta}\) is the vector of GWAS SNP effects and
\(\mathbf{R}\) a LD correlation matrix between SNPs estimated from a
reference sample. Note that the SBLUP method is equivalent to the
\texttt{LDpred-Inf} method. A detailed tutorial for LDpred2 is available
at \url{https://privefl.github.io/bigsnpr/articles/LDpred2.html}. In
this practical, we will calculate the SBLUP PGS weights ``manually'' in
R.

\textbf{Method 3: SBayes C}

The Summary-data based BayesC (SBayesC) method is an approximation of
the Bayes C method. This method is built upon a more flexible assumption
regarding the distribution of SNP effects that only allows a subset of
all \(M\) SNPs to affect the trait/disease of interest. We denote
\(\pi\) the proportion of SNPs with a non-zero effect on the trait /
disease. (S)BayesC implements a prior mixture distribution for the joint
SNP effects \(\mathbf{b} = \left(b_1,\ldots,b_M\right)\) defined as

\[ b_j \sim \pi \mathcal{N}\left(0,\sigma^2_b\right)  + (1-\pi) \delta_0\]

Note that the SBayesC method is equivalent to the \texttt{LDpred}
(non-infinitesimal) method. In this practical, we will use the GCTB
software to fit SBayes C.

\newpage

\textbf{Overview of the data}

We provide GWAS summary statistics for two (simulated) traits hereafter
denoted \texttt{Trait1} and \texttt{Trait2}. For simplicity, we will
focus summary statistics of \(M=32,260\) SNPs located on chromosome 20.
These two GWAS were conducted in \(N=348,501\) unrelated European
ancestry participants from the UK Biobank (data accessed under project
12505) using the \texttt{fastGWA} module of the \texttt{GCTA} software
(Yang et al.~2011).

We will assess prediction \(n=2504\) samples from the 1000 Genomes
Project (1KG) with different ancestries. Ancestries groups are European
(EUR; N=503), East-Asian (EAS; N=504), South-Asian (SAS; N=489), Admixed
individuals from the Americas (AMR; N=347) and African (AFR; N=661). We
provide genotypes under binary PLINK format (three files:
\texttt{1kg\_hm3.bed}, \texttt{1kg\_hm3.bim} and \texttt{1kg\_hm3.fam})
for these 2504 samples and phenotypes (files named:
\texttt{1kg.Trait1.phen} and \texttt{1kg.Trait2.phen}). For more details
about this format please visit PLINK website
(\url{https://www.cog-genomics.org/plink/1.9/formats\#bed}).

Create a folder for the practical and move in that folder

\begin{Shaded}
\begin{Highlighting}[]
\FunctionTok{mkdir}\NormalTok{ Practical6}
\BuiltInTok{cd}\NormalTok{ Practical6}
\end{Highlighting}
\end{Shaded}

Copy the data needed for the practical.

\begin{enumerate}
\def\labelenumi{\arabic{enumi})}
\tightlist
\item
  Copy the genotype files (data from 1000 Genomes Project; 1KG)
\end{enumerate}

\begin{Shaded}
\begin{Highlighting}[]
\FunctionTok{cp}\NormalTok{ /data/SISG2023M15/data/1kg\_hm3.}\PreprocessorTok{*}\NormalTok{ .}
\end{Highlighting}
\end{Shaded}

\begin{enumerate}
\def\labelenumi{\arabic{enumi})}
\setcounter{enumi}{1}
\tightlist
\item
  Copy GWAS summary statistics generated using \texttt{fastGWA}
\end{enumerate}

\begin{Shaded}
\begin{Highlighting}[]
\FunctionTok{cp}\NormalTok{ /data/SISG2023M15/data/}\PreprocessorTok{*}\NormalTok{.fastGWA .}
\end{Highlighting}
\end{Shaded}

\begin{enumerate}
\def\labelenumi{\arabic{enumi})}
\setcounter{enumi}{2}
\tightlist
\item
  Copy phenotype data for the validation sample used to assess
  prediction accuracy
\end{enumerate}

\begin{Shaded}
\begin{Highlighting}[]
\FunctionTok{cp}\NormalTok{ /data/SISG2023M15/data/}\PreprocessorTok{*}\NormalTok{.phen .}
\end{Highlighting}
\end{Shaded}

\begin{enumerate}
\def\labelenumi{\arabic{enumi})}
\setcounter{enumi}{3}
\tightlist
\item
  Copy ancestry information for the validation sample
\end{enumerate}

\begin{Shaded}
\begin{Highlighting}[]
\FunctionTok{cp}\NormalTok{ /data/SISG2023M15/data/1kg{-}sample{-}2504{-}phased.txt .}
\end{Highlighting}
\end{Shaded}

\begin{enumerate}
\def\labelenumi{\arabic{enumi})}
\setcounter{enumi}{4}
\tightlist
\item
  Copy list of 1KG samples with European ancestries
\end{enumerate}

\begin{Shaded}
\begin{Highlighting}[]
\FunctionTok{cp}\NormalTok{ /data/SISG2023M15/data/EUR.id .}
\end{Highlighting}
\end{Shaded}

\subsection{Part 1. Prediction using
C+PT.}\label{part-1.-prediction-using-cpt.}

Create a sub-folder to store all intermediate files generated for this
method.

\begin{Shaded}
\begin{Highlighting}[]
\FunctionTok{mkdir}\NormalTok{ C+PT}
\end{Highlighting}
\end{Shaded}

We can run a first example using \texttt{Trait1} with following clumping
parameters for PLINK: p-value below 5e-8 and squared correlation below
\texttt{r2\_thres=0.1} to determine independence between SNPs located
within a window \texttt{window\_kb=1000} kb (=1,000,000 base pairs). The
command lines below allow you to that

\begin{Shaded}
\begin{Highlighting}[]
\VariableTok{window\_kb}\OperatorTok{=}\NormalTok{1000 }\CommentTok{\# 1000 kb = 1 Mb window}
\VariableTok{r2\_thresh}\OperatorTok{=}\NormalTok{0.1  }\CommentTok{\# LD threshold for clumping}
\VariableTok{trait}\OperatorTok{=}\NormalTok{Trait1}
\VariableTok{pv\_thresh}\OperatorTok{=}\NormalTok{5e{-}8}

\VariableTok{PLINK}\OperatorTok{=}\NormalTok{/data/SISG2023M15/exe/plink}
\VariableTok{$\{PLINK\}} \AttributeTok{{-}{-}bfile}\NormalTok{ 1kg\_hm3 }\DataTypeTok{\textbackslash{}}
      \AttributeTok{{-}{-}keep}\NormalTok{ EUR.id }\DataTypeTok{\textbackslash{}}
      \AttributeTok{{-}{-}clump} \VariableTok{$\{trait\}}\NormalTok{.fastGWA }\DataTypeTok{\textbackslash{}}
      \AttributeTok{{-}{-}clump{-}kb} \VariableTok{$\{window\_kb\}} \DataTypeTok{\textbackslash{}}
      \AttributeTok{{-}{-}clump{-}p1} \VariableTok{$\{pv\_thresh\}} \DataTypeTok{\textbackslash{}}
      \AttributeTok{{-}{-}clump{-}p2} \VariableTok{$\{pv\_thresh\}} \DataTypeTok{\textbackslash{}}
      \AttributeTok{{-}{-}clump{-}r2} \VariableTok{$\{r2\_thresh\}} \DataTypeTok{\textbackslash{}}
      \AttributeTok{{-}{-}out}\NormalTok{ C+PT/}\VariableTok{$\{trait\}}\NormalTok{\_rsq\_}\VariableTok{$\{r2\_thresh\}}\NormalTok{\_p\_below\_}\VariableTok{$\{pv\_thresh\}}
\end{Highlighting}
\end{Shaded}

\textbf{Question 1. How many SNPs were selected using the command
above?}

We will now consider multiple significance thresholds and also analyse
the two traits. You could use the code below

\begin{Shaded}
\begin{Highlighting}[]
\ControlFlowTok{for}\NormalTok{ pv\_thresh }\KeywordTok{in}\NormalTok{ 5e{-}8 5e{-}7 5e{-}6 5e{-}5 5e{-}4 5e{-}3 5e{-}2}
\ControlFlowTok{do}
  \ControlFlowTok{for}\NormalTok{ trait }\KeywordTok{in}\NormalTok{ Trait1 Trait2}
    \ControlFlowTok{do}
    \VariableTok{$\{PLINK\}} \AttributeTok{{-}{-}bfile}\NormalTok{ 1kg\_hm3 }\DataTypeTok{\textbackslash{}}
      \AttributeTok{{-}{-}keep}\NormalTok{ EUR.id }\DataTypeTok{\textbackslash{}}
      \AttributeTok{{-}{-}clump} \VariableTok{$\{trait\}}\NormalTok{.fastGWA }\DataTypeTok{\textbackslash{}}
      \AttributeTok{{-}{-}clump{-}kb} \VariableTok{$\{window\_kb\}} \DataTypeTok{\textbackslash{}}
      \AttributeTok{{-}{-}clump{-}p1} \VariableTok{$\{pv\_thresh\}} \DataTypeTok{\textbackslash{}}
      \AttributeTok{{-}{-}clump{-}p2} \VariableTok{$\{pv\_thresh\}} \DataTypeTok{\textbackslash{}}
      \AttributeTok{{-}{-}clump{-}r2} \VariableTok{$\{r2\_thresh\}} \DataTypeTok{\textbackslash{}}
      \AttributeTok{{-}{-}out}\NormalTok{ C+PT/}\VariableTok{$\{trait\}}\NormalTok{\_rsq\_}\VariableTok{$\{r2\_thresh\}}\NormalTok{\_p\_below\_}\VariableTok{$\{pv\_thresh\}} \AttributeTok{{-}{-}silent}
  \ControlFlowTok{done}
\ControlFlowTok{done}
\end{Highlighting}
\end{Shaded}

Now we can calculate the PGS with PLINK using the \texttt{-\/-score}
command (help: https://www.cog-genomics.org/plink/1.9/score).

\begin{Shaded}
\begin{Highlighting}[]
\CommentTok{\#\# Calculate PGS for each predictor}
\ControlFlowTok{for}\NormalTok{ pv\_thresh }\KeywordTok{in}\NormalTok{ 5e{-}8 5e{-}7 5e{-}6 5e{-}5 5e{-}4 5e{-}3 5e{-}2}
  \ControlFlowTok{do}
  \ControlFlowTok{for}\NormalTok{ trait }\KeywordTok{in}\NormalTok{ Trait1 Trait2}
    \ControlFlowTok{do}
    \VariableTok{$\{PLINK\}} \AttributeTok{{-}{-}bfile}\NormalTok{ 1kg\_hm3 }\DataTypeTok{\textbackslash{}}
      \AttributeTok{{-}{-}score} \VariableTok{$\{trait\}}\NormalTok{.fastGWA 2 4 8 sum center }\DataTypeTok{\textbackslash{}}
      \AttributeTok{{-}{-}extract}\NormalTok{ C+PT/}\VariableTok{$\{trait\}}\NormalTok{\_rsq\_}\VariableTok{$\{r2\_thresh\}}\NormalTok{\_p\_below\_}\VariableTok{$\{pv\_thresh\}}\NormalTok{.clumped }\DataTypeTok{\textbackslash{}}
      \AttributeTok{{-}{-}out}\NormalTok{ C+PT/}\VariableTok{$\{trait\}}\NormalTok{\_rsq\_}\VariableTok{$\{r2\_thresh\}}\NormalTok{\_p\_below\_}\VariableTok{$\{pv\_thresh\}}\NormalTok{.pred }\AttributeTok{{-}{-}silent}
  \ControlFlowTok{done}
\ControlFlowTok{done}
\end{Highlighting}
\end{Shaded}

PGS calculated using the command above are stored in the
\texttt{*.profile} files (last column named \texttt{SCORESUM}). We can
now load them in \texttt{R} and calculate prediction accuracy in each
population.

\begin{Shaded}
\begin{Highlighting}[]
\FunctionTok{options}\NormalTok{(}\AttributeTok{stringsAsFactors =} \ConstantTok{FALSE}\NormalTok{)}
\NormalTok{r2\_thr }\OtherTok{\textless{}{-}} \FloatTok{0.1}
\NormalTok{pops   }\OtherTok{\textless{}{-}} \FunctionTok{read.table}\NormalTok{(}\StringTok{"1kg{-}sample{-}2504{-}phased.txt"}\NormalTok{,}\AttributeTok{h=}\NormalTok{T) }\DocumentationTok{\#\# This contains ancestry information}
\NormalTok{trait  }\OtherTok{\textless{}{-}} \StringTok{"Trait2"}
\NormalTok{phen   }\OtherTok{\textless{}{-}} \FunctionTok{read.table}\NormalTok{(}\FunctionTok{paste0}\NormalTok{(}\StringTok{"1kg."}\NormalTok{,trait,}\StringTok{".phen"}\NormalTok{),}\AttributeTok{h=}\NormalTok{T)[,}\SpecialCharTok{{-}}\DecValTok{1}\NormalTok{] }\DocumentationTok{\#\# Read phenotype}
\FunctionTok{colnames}\NormalTok{(phen) }\OtherTok{\textless{}{-}} \FunctionTok{c}\NormalTok{(}\StringTok{"IID"}\NormalTok{,trait)}

\DocumentationTok{\#\# Merge with all PGS}
\NormalTok{Threshs }\OtherTok{\textless{}{-}} \DecValTok{8}\SpecialCharTok{:}\DecValTok{2}
\ControlFlowTok{for}\NormalTok{(thresh }\ControlFlowTok{in}\NormalTok{ Threshs)\{}
\NormalTok{  filename }\OtherTok{\textless{}{-}} \FunctionTok{paste0}\NormalTok{(}\StringTok{"C+PT/"}\NormalTok{,trait,}\StringTok{"\_rsq\_"}\NormalTok{,r2\_thr,}\StringTok{"\_p\_below\_5e{-}"}\NormalTok{,thresh,}\StringTok{".pred.profile"}\NormalTok{)}
\NormalTok{  tmp }\OtherTok{\textless{}{-}} \FunctionTok{read.table}\NormalTok{(filename,}\AttributeTok{h=}\NormalTok{T)[,}\FunctionTok{c}\NormalTok{(}\DecValTok{2}\NormalTok{,}\DecValTok{6}\NormalTok{)]}
  \FunctionTok{colnames}\NormalTok{(tmp)[}\DecValTok{2}\NormalTok{] }\OtherTok{\textless{}{-}} \FunctionTok{paste0}\NormalTok{(}\StringTok{"P"}\NormalTok{,thresh)}
\NormalTok{  phen }\OtherTok{\textless{}{-}} \FunctionTok{merge}\NormalTok{(phen,tmp,}\AttributeTok{by=}\StringTok{"IID"}\NormalTok{)}
\NormalTok{\}}

\DocumentationTok{\#\# This is the merged data containing all PGS}
\DocumentationTok{\#\# P8 denotes the PGS calculated using SNPs with p{-}value \textless{}5e{-}8}
\NormalTok{phen }\OtherTok{\textless{}{-}} \FunctionTok{merge}\NormalTok{(phen,pops,}\AttributeTok{by.x=}\StringTok{"IID"}\NormalTok{,}\AttributeTok{by.y=}\StringTok{"sample"}\NormalTok{)}

\DocumentationTok{\#\# Look at prediction accuracy}
\CommentTok{\# across the sample}
\FunctionTok{summary}\NormalTok{(}\FunctionTok{lm}\NormalTok{(}\FunctionTok{paste0}\NormalTok{(trait,}\StringTok{"\textasciitilde{}P8"}\NormalTok{),phen)) }

\DocumentationTok{\#\# within European ancestry individuals}
\FunctionTok{summary}\NormalTok{(}\FunctionTok{lm}\NormalTok{(}\FunctionTok{paste0}\NormalTok{(trait,}\StringTok{"\textasciitilde{}P8"}\NormalTok{),phen[}\FunctionTok{which}\NormalTok{(phen}\SpecialCharTok{$}\NormalTok{super\_pop}\SpecialCharTok{==}\StringTok{"EUR"}\NormalTok{),])) }

\DocumentationTok{\#\# within African  ancestry individuals}
\FunctionTok{summary}\NormalTok{(}\FunctionTok{lm}\NormalTok{(}\FunctionTok{paste0}\NormalTok{(trait,}\StringTok{"\textasciitilde{}P8"}\NormalTok{),phen[}\FunctionTok{which}\NormalTok{(phen}\SpecialCharTok{$}\NormalTok{super\_pop}\SpecialCharTok{==}\StringTok{"AFR"}\NormalTok{),])) }

\DocumentationTok{\#\# Select best threshold}
\DocumentationTok{\#\# Across all samples}
\NormalTok{RsqOverall }\OtherTok{\textless{}{-}} \FunctionTok{rep}\NormalTok{(}\ConstantTok{NA}\NormalTok{,}\FunctionTok{length}\NormalTok{(Threshs))}
\FunctionTok{names}\NormalTok{(RsqOverall) }\OtherTok{\textless{}{-}} \FunctionTok{paste0}\NormalTok{(}\StringTok{"P"}\NormalTok{,Threshs)}
\ControlFlowTok{for}\NormalTok{(thresh }\ControlFlowTok{in}\NormalTok{ Threshs)\{}
\NormalTok{  idThresh }\OtherTok{\textless{}{-}} \FunctionTok{paste0}\NormalTok{(}\StringTok{"P"}\NormalTok{,thresh)}
\NormalTok{  RsqOverall[idThresh] }\OtherTok{\textless{}{-}} \FunctionTok{cor}\NormalTok{(phen[,trait],phen[,idThresh])}\SpecialCharTok{\^{}}\DecValTok{2}
\NormalTok{\}}

\DocumentationTok{\#\# Within each ancestry group}
\NormalTok{ancestries }\OtherTok{\textless{}{-}} \FunctionTok{c}\NormalTok{(}\StringTok{"EUR"}\NormalTok{,}\StringTok{"SAS"}\NormalTok{,}\StringTok{"EAS"}\NormalTok{,}\StringTok{"AMR"}\NormalTok{,}\StringTok{"AFR"}\NormalTok{)}
\NormalTok{RsqPerAnc  }\OtherTok{\textless{}{-}} \FunctionTok{matrix}\NormalTok{(}\ConstantTok{NA}\NormalTok{,}\AttributeTok{nrow=}\FunctionTok{length}\NormalTok{(ancestries),}\AttributeTok{ncol=}\FunctionTok{length}\NormalTok{(Threshs))}
\FunctionTok{colnames}\NormalTok{(RsqPerAnc) }\OtherTok{\textless{}{-}} \FunctionTok{paste0}\NormalTok{(}\StringTok{"P"}\NormalTok{,Threshs)}
\FunctionTok{rownames}\NormalTok{(RsqPerAnc) }\OtherTok{\textless{}{-}}\NormalTok{ ancestries}
\ControlFlowTok{for}\NormalTok{(thresh }\ControlFlowTok{in}\NormalTok{ Threshs)\{}
\NormalTok{  idThresh }\OtherTok{\textless{}{-}} \FunctionTok{paste0}\NormalTok{(}\StringTok{"P"}\NormalTok{,thresh)}
  \ControlFlowTok{for}\NormalTok{(ancestry }\ControlFlowTok{in}\NormalTok{ ancestries)\{}
\NormalTok{    idAnc }\OtherTok{\textless{}{-}} \FunctionTok{which}\NormalTok{(phen[,}\StringTok{"super\_pop"}\NormalTok{]}\SpecialCharTok{==}\NormalTok{ancestry)}
\NormalTok{    RsqPerAnc[ancestry,idThresh] }\OtherTok{\textless{}{-}} \FunctionTok{cor}\NormalTok{(phen[idAnc,trait],phen[idAnc,idThresh])}\SpecialCharTok{\^{}}\DecValTok{2}
\NormalTok{  \}}
\NormalTok{\}}
\NormalTok{RsqPerAnc}
\end{Highlighting}
\end{Shaded}

\textbf{Question 2. }

\textbf{What is the best prediction accuracy in EUR individuals across
all significance thresholds? Is the best accuracy reached using the same
significance threshold consistently across all ancestry groups? In which
ancestry group is the prediction higher? Do you see any difference
between \texttt{Trait1} and \texttt{Trait2}?}

If you have time (\textbf{at the end of the Practical}) and answered the
question above then you could split the EUR sample in 2 (random)
subsets, select the threshold that yields the larger accuracy in
subset1, then re-evaluate the prediction accuracy of the corresponding
PGS in subset 2. Compare the resulting prediction accuracy with your
answer to the first part of \textbf{Question 2} above. This process is
called cross-validation.

\subsection{Part 2. Prediction using
SBLUP.}\label{part-2.-prediction-using-sblup.}

Create a folder to store all intermediate files generated for this
method.

(Back to the terminal!)

\begin{Shaded}
\begin{Highlighting}[]
\FunctionTok{mkdir}\NormalTok{ SBLUP}
\end{Highlighting}
\end{Shaded}

The SNP-based heritability of \texttt{Trait1} and \texttt{Trait2} is
\(h^2_{SNP} = 0.2\). As shown above, the SBLUP uses a LD reference. Here
we provide LD data calculated with PLINK in \(N=348,501\) unrelated
European ancestry participants from the UK Biobank using the following
command.

\begin{Shaded}
\begin{Highlighting}[]
\CommentTok{\#\# Calculate PGS for each predictor}
\VariableTok{$\{PLINK\}} \DataTypeTok{\textbackslash{}}
  \AttributeTok{{-}{-}bfile}\NormalTok{ UKBu\_chrom20 }\DataTypeTok{\textbackslash{}}
  \AttributeTok{{-}{-}chr}\NormalTok{ 20 }\DataTypeTok{\textbackslash{}}
  \AttributeTok{{-}{-}ld{-}window}\NormalTok{ 999999999 }\DataTypeTok{\textbackslash{}}
  \AttributeTok{{-}{-}ld{-}window{-}kb}\NormalTok{ 1000 }\DataTypeTok{\textbackslash{}}
  \AttributeTok{{-}{-}ld{-}window{-}r2}\NormalTok{ 0.00 }\DataTypeTok{\textbackslash{}}
  \AttributeTok{{-}{-}r} \StringTok{\textquotesingle{}yes{-}really\textquotesingle{}} \StringTok{\textquotesingle{}gz\textquotesingle{}} \DataTypeTok{\textbackslash{}}
  \AttributeTok{{-}{-}out}\NormalTok{ plinkLDMat\_chrom20}
\end{Highlighting}
\end{Shaded}

{\textbf{We only show this command out of general interest. For the sake
of time, we do not advise running this command during the practical but
you could try it later at home. Note that the data named
\texttt{UKBu\_chrom20} has not been provided but you could replace it
another dataset, e.g., with \texttt{1kg\_hm3}.}}

You will see below an R script to run SBLUP. This code is provided in an
external \texttt{R} script named \texttt{sblup.R}. You can run it using
the following command

\begin{Shaded}
\begin{Highlighting}[]
\CommentTok{\#\# Calculate PGS for each predictor}
\ExtensionTok{Rscript}\NormalTok{ /data/SISG2023M15/data/sblup.R Trait1}
\end{Highlighting}
\end{Shaded}

The command above will generate a file named
\texttt{SBLUP/Trait1.sblupInR.res}.

\textbf{Question 3. Using the examples above,1) calculate the PGS using
SNP effects from SBLUP (for \texttt{Trait1} and \texttt{Trait2}); 2)
assess prediction accuracy in the test sample 3) compare your results
with the prediction accuracy in EUR and non-EUR individuals obtained
using the C+PT method}.

It is expected that you try to answer these questions using the R
functions introduced above. The below code is a partial solution to the
above questions and uses some key functions that you may want to use to
answer the above questions.

\begin{Shaded}
\begin{Highlighting}[]
\FunctionTok{options}\NormalTok{(}\AttributeTok{stringsAsFactors =} \ConstantTok{FALSE}\NormalTok{)}
\NormalTok{pops   }\OtherTok{\textless{}{-}} \FunctionTok{read.table}\NormalTok{(}\StringTok{"1kg{-}sample{-}2504{-}phased.txt"}\NormalTok{,}\AttributeTok{h=}\NormalTok{T)}
\NormalTok{trait  }\OtherTok{\textless{}{-}} \StringTok{"Trait2"}
\NormalTok{phen   }\OtherTok{\textless{}{-}} \FunctionTok{read.table}\NormalTok{(}\FunctionTok{paste0}\NormalTok{(}\StringTok{"1kg."}\NormalTok{,trait,}\StringTok{".phen"}\NormalTok{),}\AttributeTok{h=}\NormalTok{T)[,}\SpecialCharTok{{-}}\DecValTok{1}\NormalTok{] }\DocumentationTok{\#\# Read phenotype}
\FunctionTok{colnames}\NormalTok{(phen) }\OtherTok{\textless{}{-}} \FunctionTok{c}\NormalTok{(}\StringTok{"IID"}\NormalTok{,trait)}
\NormalTok{sblup  }\OtherTok{\textless{}{-}} \FunctionTok{read.table}\NormalTok{(}\FunctionTok{paste0}\NormalTok{(}\StringTok{"SBLUP/"}\NormalTok{,trait,}\StringTok{".pred.R.profile"}\NormalTok{),}\AttributeTok{h=}\NormalTok{T)[,}\FunctionTok{c}\NormalTok{(}\DecValTok{2}\NormalTok{,}\DecValTok{6}\NormalTok{)]}

\FunctionTok{colnames}\NormalTok{(sblup)[}\DecValTok{2}\NormalTok{] }\OtherTok{\textless{}{-}} \StringTok{"SBLUP"}
\NormalTok{phen }\OtherTok{\textless{}{-}} \FunctionTok{merge}\NormalTok{(phen,sblup,}\AttributeTok{by=}\StringTok{"IID"}\NormalTok{)}
\NormalTok{phen }\OtherTok{\textless{}{-}} \FunctionTok{merge}\NormalTok{(phen,pops,}\AttributeTok{by.x=}\StringTok{"IID"}\NormalTok{,}\AttributeTok{by.y=}\StringTok{"sample"}\NormalTok{)}

\DocumentationTok{\#\# Look at prediction accuracy}
\CommentTok{\# across the entire sample}
\FunctionTok{summary}\NormalTok{(}\FunctionTok{lm}\NormalTok{(}\FunctionTok{paste0}\NormalTok{(trait,}\StringTok{"\textasciitilde{}SBLUP"}\NormalTok{),phen)) }

\DocumentationTok{\#\# within European ancestry individuals}
\FunctionTok{summary}\NormalTok{(}\FunctionTok{lm}\NormalTok{(}\FunctionTok{paste0}\NormalTok{(trait,}\StringTok{"\textasciitilde{}SBLUP"}\NormalTok{),phen[}\FunctionTok{which}\NormalTok{(phen}\SpecialCharTok{$}\NormalTok{super\_pop}\SpecialCharTok{==}\StringTok{"EUR"}\NormalTok{),])) }

\DocumentationTok{\#\# within African  ancestry individuals}
\FunctionTok{summary}\NormalTok{(}\FunctionTok{lm}\NormalTok{(}\FunctionTok{paste0}\NormalTok{(trait,}\StringTok{"\textasciitilde{}SBLUP"}\NormalTok{),phen[}\FunctionTok{which}\NormalTok{(phen}\SpecialCharTok{$}\NormalTok{super\_pop}\SpecialCharTok{==}\StringTok{"AFR"}\NormalTok{),])) }

\DocumentationTok{\#\# Within each ancestry group}
\NormalTok{ancestries }\OtherTok{\textless{}{-}} \FunctionTok{c}\NormalTok{(}\StringTok{"EUR"}\NormalTok{,}\StringTok{"SAS"}\NormalTok{,}\StringTok{"EAS"}\NormalTok{,}\StringTok{"AMR"}\NormalTok{,}\StringTok{"AFR"}\NormalTok{)}
\NormalTok{RsqPerAnc  }\OtherTok{\textless{}{-}} \FunctionTok{rep}\NormalTok{(}\ConstantTok{NA}\NormalTok{,}\FunctionTok{length}\NormalTok{(ancestries))}
\FunctionTok{names}\NormalTok{(RsqPerAnc) }\OtherTok{\textless{}{-}}\NormalTok{ ancestries}
\ControlFlowTok{for}\NormalTok{(ancestry }\ControlFlowTok{in}\NormalTok{ ancestries)\{}
\NormalTok{  idAnc }\OtherTok{\textless{}{-}} \FunctionTok{which}\NormalTok{(phen[,}\StringTok{"super\_pop"}\NormalTok{]}\SpecialCharTok{==}\NormalTok{ancestry)}
\NormalTok{  RsqPerAnc[ancestry] }\OtherTok{\textless{}{-}} \FunctionTok{cor}\NormalTok{(phen[idAnc,trait],phen[idAnc,}\StringTok{"SBLUP"}\NormalTok{])}\SpecialCharTok{\^{}}\DecValTok{2}
\NormalTok{\}}
\end{Highlighting}
\end{Shaded}

It is also possible to run SBLUP using \texttt{GCTA}. This will require
a different formatting of GWAS summary statistics, which we provide
(files named \texttt{Trait1.ma} and \texttt{Trait2.ma}). Note that the
same format can also be used with \texttt{GCTB} to run SBayesC (see
\textbf{Part 3} below). The implementation of SBLUP in \texttt{GCTA}
requires genotype data as input (and not a pre-calculated LD correlation
matrix) as well as shrinkage parameter
\(\lambda = M_{SNPs}\left(\frac{1}{h^2_{SNP}}-1\right)\). An example
command is given below

\begin{Shaded}
\begin{Highlighting}[]
\VariableTok{GCTA}\OperatorTok{=}\NormalTok{/data/SISG2023M15/exe/gcta{-}1.94.1}
\VariableTok{lambda}\OperatorTok{=}\NormalTok{135628}
\ControlFlowTok{for}\NormalTok{ trait }\KeywordTok{in}\NormalTok{ Trait1 Trait2}
  \ControlFlowTok{do}
  \VariableTok{$\{GCTA\}} \DataTypeTok{\textbackslash{}}
  \AttributeTok{{-}{-}bfile}\NormalTok{ 1kg\_hm3 }\DataTypeTok{\textbackslash{}}
  \AttributeTok{{-}{-}chr}\NormalTok{ 20 }\DataTypeTok{\textbackslash{}}
  \AttributeTok{{-}{-}cojo{-}file}\NormalTok{ /data/SISG2023M15/data/}\VariableTok{$\{trait\}}\NormalTok{.ma }\DataTypeTok{\textbackslash{}}
  \AttributeTok{{-}{-}cojo{-}sblup} \VariableTok{$\{lambda\}} \DataTypeTok{\textbackslash{}}
  \AttributeTok{{-}{-}cojo{-}wind}\NormalTok{ 1000 }\DataTypeTok{\textbackslash{}}
  \AttributeTok{{-}{-}thread{-}num}\NormalTok{ 20 }\DataTypeTok{\textbackslash{}}
  \AttributeTok{{-}{-}out}\NormalTok{ SBLUP/}\VariableTok{$\{trait\}}
\ControlFlowTok{done}
\end{Highlighting}
\end{Shaded}

This command will take approximately 5 min to run. In the interest of
time, we recommend not running it and moving to \textbf{Part 3} instead.
More details can be found here:
https://yanglab.westlake.edu.cn/software/gcta/\#SBLUP.

\subsection{Part 3. Prediction using SBayes
C}\label{part-3.-prediction-using-sbayes-c}

Create a folder to store all intermediate files generated for this
method.

\begin{Shaded}
\begin{Highlighting}[]
\FunctionTok{mkdir}\NormalTok{ SBC}
\end{Highlighting}
\end{Shaded}

SBayesC requires a different format for the LD matrix. We provide a LD
matrix directly calculated with GCTB using a command like this

\begin{Shaded}
\begin{Highlighting}[]
\VariableTok{$\{GCTB\}} \AttributeTok{{-}{-}bfile}\NormalTok{ sampleForLD }\AttributeTok{{-}{-}make{-}full{-}ldm} \AttributeTok{{-}{-}snp}\NormalTok{ 1{-}5000 }\AttributeTok{{-}{-}out}\NormalTok{ SBC/BLOCK1618.CHROM20}
\end{Highlighting}
\end{Shaded}

\texttt{sampleForLD} is a subset of 348,501 unrelated EUR participants
of the UKB. LD correlations on chromosome 20 were calculated within 38
non-overlapping LD blocks identified in an EUR sample (see:
http://bitbucket.org/nygcresearch/ldetect-data). Pre-calculated LD
matrices can be downloaded from the \texttt{GCTB} website. We provide
these 38 LD submatrices in the folder named \texttt{ldm} and paths for
each file are listed in \texttt{mldm.txt} (note the similarity with
\texttt{GCTA}: in \texttt{GCTB} \texttt{-\/-mldm} indicates that
multiple LD matrices will be used, while in \texttt{GCTA}
\texttt{-\/-mgrm} indicates that multiple GRMs will be used).

Set the path to GCTB

\begin{Shaded}
\begin{Highlighting}[]
\VariableTok{GCTB}\OperatorTok{=}\NormalTok{/data/SISG2023M15/exe/gctb}
\end{Highlighting}
\end{Shaded}

Now, to run SBayes C, you could run the following command.

\begin{Shaded}
\begin{Highlighting}[]
\VariableTok{mldm}\OperatorTok{=}\NormalTok{/data/SISG2023M15/data/mldm.txt}
\VariableTok{trait}\OperatorTok{=}\NormalTok{Trait1}

\VariableTok{$\{GCTB\}} \AttributeTok{{-}{-}sbayes}\NormalTok{ C }\AttributeTok{{-}{-}mldm} \VariableTok{$\{mldm\}} \DataTypeTok{\textbackslash{}}
\NormalTok{{-}{-}gwas{-}summary /data/SISG2023M15/data/}\VariableTok{$\{trait\}}\NormalTok{.ma }\DataTypeTok{\textbackslash{}}
\NormalTok{{-}{-}pi 0.0001 }\AttributeTok{{-}{-}hsq}\NormalTok{ 0.001 }\DataTypeTok{\textbackslash{}}
\NormalTok{{-}{-}chain{-}length 10000 }\DataTypeTok{\textbackslash{}}
\NormalTok{{-}{-}burn{-}in 5000 }\DataTypeTok{\textbackslash{}}
\NormalTok{{-}{-}no{-}mcmc{-}bin }\DataTypeTok{\textbackslash{}}
\NormalTok{{-}{-}robust }\DataTypeTok{\textbackslash{}}
\NormalTok{{-}{-}out{-}freq 100 }\DataTypeTok{\textbackslash{}}
\NormalTok{{-}{-}thin 10 }\DataTypeTok{\textbackslash{}}
\NormalTok{{-}{-}out SBC/}\VariableTok{$\{trait\}}
\end{Highlighting}
\end{Shaded}

This command will run 10,000 MCMC iterations (\texttt{-\/-chain-length})
and output results every 100 iterations (\texttt{-\/-out-freq}). Final
statistics for parameters and prediction accuracy will be calculated
using 1/10 iterations (\texttt{-\/-thin}) after a burn-in of 5000
iterations (\texttt{-\/-burn-in}).

\textbf{Question 4. What is the heritability of \texttt{Trait1}? What is
the proportion/expected number of SNPs with a non-zero effect on the
trait? What is the genetic and residual variances?}.

The command above will generate a file named \texttt{SBC/Trait1.snpRes}.
We can use that file to calculated a PGS using the following command

\begin{Shaded}
\begin{Highlighting}[]
\VariableTok{trait}\OperatorTok{=}\NormalTok{Trait1}
\VariableTok{$\{PLINK\}} \AttributeTok{{-}{-}bfile}\NormalTok{ 1kg\_hm3 }\DataTypeTok{\textbackslash{}}
\NormalTok{{-}{-}score SBC/}\VariableTok{$\{trait\}}\NormalTok{.snpRes 2 5 8 sum center }\DataTypeTok{\textbackslash{}}
\NormalTok{{-}{-}out SBC/}\VariableTok{$\{trait\}}\NormalTok{.pred}
\end{Highlighting}
\end{Shaded}

Assess the prediction accuracy of the SBayesC PGS in \texttt{R} as done
previously

\begin{Shaded}
\begin{Highlighting}[]
\ExtensionTok{options}\ErrorTok{(}\ExtensionTok{stringsAsFactors}\NormalTok{ = FALSE}\KeywordTok{)}
\ExtensionTok{pops}   \OperatorTok{\textless{}}\NormalTok{{-} read.table}\ErrorTok{(}\StringTok{"1kg{-}sample{-}2504{-}phased.txt"}\ExtensionTok{,h=T}\KeywordTok{)}
\ExtensionTok{trait}  \OperatorTok{\textless{}}\NormalTok{{-} }\StringTok{"Trait1"}
\ExtensionTok{phen}   \OperatorTok{\textless{}}\NormalTok{{-} read.table}\ErrorTok{(}\ExtensionTok{paste0}\ErrorTok{(}\StringTok{"1kg."}\ExtensionTok{,trait,}\StringTok{".phen"}\KeywordTok{)}\ExtensionTok{,h=T}\KeywordTok{)}\ExtensionTok{[,{-}1]} \CommentTok{\#\# Read phenotype}
\ExtensionTok{colnames}\ErrorTok{(}\ExtensionTok{phen}\KeywordTok{)} \OperatorTok{\textless{}}\NormalTok{{-} }\ExtensionTok{c}\ErrorTok{(}\StringTok{"IID"}\ExtensionTok{,trait}\KeywordTok{)}
\ExtensionTok{sblup}  \OperatorTok{\textless{}}\NormalTok{{-} read.table}\ErrorTok{(}\ExtensionTok{paste0}\ErrorTok{(}\StringTok{"SBC/"}\ExtensionTok{,trait,}\StringTok{".pred.profile"}\KeywordTok{)}\ExtensionTok{,h=T}\KeywordTok{)}\ExtensionTok{[,c}\ErrorTok{(}\ExtensionTok{2,6}\KeywordTok{)}\ExtensionTok{]}
\ExtensionTok{colnames}\ErrorTok{(}\ExtensionTok{sblup}\KeywordTok{)}\ExtensionTok{[2]} \OperatorTok{\textless{}}\NormalTok{{-} }\StringTok{"SBC"}
\ExtensionTok{phen} \OperatorTok{\textless{}}\NormalTok{{-} merge}\ErrorTok{(}\ExtensionTok{phen,sblup,by=}\StringTok{"IID"}\KeywordTok{)}
\ExtensionTok{phen} \OperatorTok{\textless{}}\NormalTok{{-} merge}\ErrorTok{(}\ExtensionTok{phen,pops,by.x=}\StringTok{"IID"}\ExtensionTok{,by.y=}\StringTok{"sample"}\KeywordTok{)}

\CommentTok{\#\# Look at prediction accuracy}
\CommentTok{\# across the sample}
\ExtensionTok{summary}\ErrorTok{(}\ExtensionTok{lm}\ErrorTok{(}\ExtensionTok{paste0}\ErrorTok{(}\ExtensionTok{trait,}\StringTok{"\textasciitilde{}SBC"}\KeywordTok{)}\ExtensionTok{,phen}\KeywordTok{))} 

\CommentTok{\#\# within European ancestry individuals}
\ExtensionTok{summary}\ErrorTok{(}\ExtensionTok{lm}\ErrorTok{(}\ExtensionTok{paste0}\ErrorTok{(}\ExtensionTok{trait,}\StringTok{"\textasciitilde{}SBC"}\KeywordTok{)}\ExtensionTok{,phen[which}\ErrorTok{(}\ExtensionTok{phen}\VariableTok{$super\_pop}\ExtensionTok{==}\StringTok{"EUR"}\KeywordTok{)}\ExtensionTok{,]}\KeywordTok{))} 

\CommentTok{\#\# within African  ancestry individuals}
\ExtensionTok{summary}\ErrorTok{(}\ExtensionTok{lm}\ErrorTok{(}\ExtensionTok{paste0}\ErrorTok{(}\ExtensionTok{trait,}\StringTok{"\textasciitilde{}SBC"}\KeywordTok{)}\ExtensionTok{,phen[which}\ErrorTok{(}\ExtensionTok{phen}\VariableTok{$super\_pop}\ExtensionTok{==}\StringTok{"AFR"}\KeywordTok{)}\ExtensionTok{,]}\KeywordTok{))} 

\CommentTok{\#\# Within each ancestry group}
\ExtensionTok{ancestries} \OperatorTok{\textless{}}\NormalTok{{-} c}\ErrorTok{(}\StringTok{"EUR"}\ExtensionTok{,}\StringTok{"SAS"}\ExtensionTok{,}\StringTok{"EAS"}\ExtensionTok{,}\StringTok{"AMR"}\ExtensionTok{,}\StringTok{"AFR"}\KeywordTok{)}
\ExtensionTok{RsqPerAnc}  \OperatorTok{\textless{}}\NormalTok{{-} rep}\ErrorTok{(}\ExtensionTok{NA,length}\ErrorTok{(}\ExtensionTok{ancestries}\KeywordTok{))}
\ExtensionTok{names}\ErrorTok{(}\ExtensionTok{RsqPerAnc}\KeywordTok{)} \OperatorTok{\textless{}}\NormalTok{{-} }\ExtensionTok{ancestries}
\ControlFlowTok{for}\KeywordTok{(}\ExtensionTok{ancestry}\NormalTok{ in ancestries}\KeywordTok{)\{}
  \ExtensionTok{idAnc} \OperatorTok{\textless{}}\NormalTok{{-} which}\ErrorTok{(}\VariableTok{phen}\OperatorTok{[}\NormalTok{,}\StringTok{"super\_pop"}\OperatorTok{]=}\NormalTok{=ancestry}\KeywordTok{)}
  \VariableTok{RsqPerAnc}\OperatorTok{[}\NormalTok{ancestry}\OperatorTok{]} \OperatorTok{\textless{}}\NormalTok{{-} }\ExtensionTok{cor}\ErrorTok{(}\VariableTok{phen}\OperatorTok{[}\NormalTok{idAnc,trait}\OperatorTok{]}\ExtensionTok{,phen[idAnc,}\StringTok{"SBC"}\ExtensionTok{]}\KeywordTok{)}\ExtensionTok{\^{}2}
\KeywordTok{\}}
\end{Highlighting}
\end{Shaded}





\end{document}
